\chapter{Hello World}
\label{chap:fpgasteel-hello-world}

\section{About this project}
\label{sec:fpgasteel-hello-about}

|hello_word| (under |sw/hello_word|) is the first application in \sysname{FPGAsteel}, and it sets the pattern every later example follows: hardware is wrapped in a small set of classes, each initialized with |init()| before use and torn down with |deInit()| afterwards. It:

\begin{itemize}
    \item brings up the interrupt controller;
    \item brings up \term{UART0};
    \item prints ``Hello, World!'' followed by ``GoodBye'';
    \item spins forever.
\end{itemize}

The purpose of this program is to check that the whole chain works end to end: |make| producing an |.elf| with the \term{ARM} toolchain, \term{OpenOCD} loading it over \term{JTAG}, and the \term{HPS} running it with interrupts enabled and talking over the serial console. Everything built afterwards assumes this chain works.

\section{Prerequisites}
\label{sec:fpgasteel-hello-prerequisites}

The toolchain must already be set up as described in Chapter~\ref{chap:fpgasteel-board-bringup}, and the \term{SD card} prepared accordingly:

\begin{itemize}
    \item the bootloader (\term{SPL} + \term{U-Boot}) written to the raw \term{A2} partition;
    \item the \term{FPGA} bitstream |soc_system.rbf| copied to the boot partition, with \term{MSEL[4:0]} on \term{SW10} set to 00000 so the \term{FPGA} is configured by the \term{HPS} via the \term{FPGA Manager};
    \item |boot.script.jtag_debug| compiled to |u-boot.scr| and copied to the boot partition, not |multiboot.script|.
\end{itemize}

If the \term{SD card} was already used with a different boot script (e.g.\ |multiboot.script|), don't just swap the file: wipe the card and redo the whole \term{SD card} preparation from scratch (Board Bring-up, Section~\ref{sec:fpgasteel-preparing-sd-card}), this time compiling |boot.script.jtag_debug| in Section~\ref{sec:fpgasteel-hello-jtag-debug}. This guarantees the board ends up at the ``Ready for JTAG'' prompt this project needs.

Two tools are needed on top of the ones from Board Bring-up, to debug over \term{JTAG} from an \term{IDE}:

\begin{itemize}
    \item \term{Eclipse IDE for Embedded C/C++}
    \item \term{OpenOCD} (\term{xPack})
\end{itemize}

See Paragraph~\ref{subsec:fpgasteel-install-paths} for versions and install paths.

\section{Code structure}
\label{sec:fpgasteel-hello-code-structure}

The project defines two main classes, |System| and |Serial|, tasked respectively with managing the interrupt controller and \term{UART0}. Since each one is tied to a single, specific piece of hardware (there is exactly one interrupt controller and one \term{UART0}), both are implemented as singletons: the caller obtains the instance with |getInstance()|, then calls |init()|, which performs the actual hardware initialization. If |init()| fails, both throw a matching exception (|ExceptionSystem| for |System|, |ExceptionSerial| for |Serial|), carrying the name of the failing \term{hwlib} call and its \term{ALT\_STATUS\_CODE}.

\subsection{System}
\label{subsec:fpgasteel-hello-system}

|System| (|src/System.{h,cpp}|) manages the interrupt controller (\term{GIC}):

\begin{itemize}
    \item |getInstance()|: returns the singleton instance.
    \item |init()|: initializes and enables the interrupt controller, in order: |alt_int_global_init()|, |alt_int_cpu_init()|, |alt_int_cpu_enable()|, |alt_int_global_enable()|.
    \item |deInit()|: disables and uninitializes it, in the reverse order.
\end{itemize}

Both |init()| and |deInit()| throw |ExceptionSystem| if any of the underlying calls fails.

\subsection{Serial}
\label{subsec:fpgasteel-hello-serial}

|Serial| (|src/Serial.{h,cpp}|) manages \term{UART0} in interrupt mode, 8-N-1 at 115200 baud (constants in |Config.h|):

\begin{itemize}
    \item |getInstance()|: returns the singleton instance.
    \item |init()|: configures \term{UART0} (data bits, parity, stop bits, baud rate), enables its \term{FIFO}s, registers the \term{UART0} interrupt handler and enables the \term{RX} and line interrupts.
    \item |deInit()|: disables \term{UART0}, unregisters the interrupt handler and discards any buffered data.
    \item |write(data, size)|: queues bytes for transmission and enables the \term{TX} interrupt.
    \item |flush()|: blocks until the \term{TX} buffer has drained and the \term{UART} has finished shifting out the last byte.
    \item |printf(format, ...)|: formats into a fixed-size stack buffer, then calls |write()|.
    \item |read(data, size)|: copies up to |size| received bytes without blocking.
    \item |available()|: number of received bytes waiting to be read.
    \item |isTxFull()| / |isReady()|: status queries.
\end{itemize}

|init()|, |deInit()|, |flush()| and |write()| throw |ExceptionSerial| if the underlying \term{hwlib} call fails.

Internally, |Serial| uses two instances of |StaticRingBuffer<char, N>| (|src/StaticRingBuffer.{h,cpp}|), one for \term{TX} and one for \term{RX}: a fixed-capacity, single-producer/single-consumer circular buffer with no dynamic allocation. |push()| throws |ExceptionRingBuffer| if the buffer is full, |pop()| throws it if the buffer is empty. This is what lets the \term{UART} interrupt handler (producer for \term{RX}, consumer for \term{TX}) and the main loop (producer for \term{TX}, consumer for \term{RX}) exchange bytes safely without a lock.

\subsection{hwlib}
\label{subsec:fpgasteel-hello-hwlib}

|System| and |Serial| sit on top of \term{hwlib}, \term{Altera}/\term{Intel}'s low-level driver library for the \term{HPS} peripherals (interrupt controller, \term{UART}, timers, \term{DMA}, etc.), pulled in as the |repos/intel-socfpga-hwlib| \term{git} submodule. Rather than depending on a prebuilt/system-wide copy, |Makefile.hwlib| compiles only what this project needs directly from the submodule's |soc_cv_av| (\term{Cyclone V}) sources into a small static library artifact, |libhwlib_socfpga.a|.

\section{How it works}
\label{sec:fpgasteel-hello-how-it-works}

The following snippet, from |main.cpp|, shows \term{System} and \term{Serial} used together, following the |init()|/|deInit()| pattern from Section~\ref{sec:fpgasteel-hello-code-structure}:

\begin{lstlisting}
System::getInstance().init();   // interrupt controller
Serial::getInstance().init();   // UART0
Serial::getInstance().write("Hello, World!\r\n", ...);
Serial::getInstance().flush();
printf("GoodBye\n");
\end{lstlisting}

|ExceptionSystem|, |ExceptionSerial| and |ExceptionRingBuffer| are caught around this block and reported back over |Serial|.

\section{Building}
\label{sec:fpgasteel-hello-building}

The project is built with a |Makefile|, from the command line. \term{Visual Studio Code} is used only as a plain text editor, not as an \term{IDE}. \term{Eclipse}, in turn, is used only for debugging: it is never used to compile the project, nor for any other \term{IDE}-specific feature.

\begin{shellcode}
cd sw/hello_word
make clean   # removes build/ and bin/
make
# -> bin/hello_word.elf, bin/hello_word.bin
\end{shellcode}

\section{Connecting the board}
\label{sec:fpgasteel-hello-connecting}

Before starting a debug session, connect the board to the host \term{PC} with two \term{USB} cables:

\begin{itemize}
    \item a \term{USB Type-B} cable, for the \term{USB Blaster II} (\term{JTAG});
    \item a \term{micro-USB} cable, for the \term{UART} terminal exposed by the \term{HPS} at 115200 baud, 8-N-1, no flow control.
\end{itemize}

On \term{Linux}, |picocom| is a lightweight terminal emulator for the latter: find the device with \texttt{dmesg \textbar{} grep tty} (typically |/dev/ttyUSB0|), then connect with:

\begin{shellcode}
picocom -b 115200 /dev/ttyUSB0
\end{shellcode}

\section{Debugging with Eclipse}
\label{sec:fpgasteel-hello-jtag-debug}

With the board wired up and sitting at ``Ready for JTAG'', the rest happens inside \term{Eclipse}: import the project, tell it how to reach the board over \term{JTAG}, then start the session.

\subsection{Import the project}
\label{subsec:fpgasteel-hello-import}

\begin{enumerate}
    \item |File > Import... > General > Existing Projects into Workspace|.
    \item Select root directory and pick the |sw/hello_word| folder and set up for the \term{Arm Cross GCC} toolchain.
    \item |Finish|.
\end{enumerate}

The imported project is not built from \term{Eclipse}; building stays |make|'s job.

\subsection{Create the Debug Configuration}
\label{subsec:fpgasteel-hello-debug-config}

\term{Eclipse} still doesn't know how to reach the board: a \term{Debug Configuration} tells it which \term{OpenOCD} config files to use and how to attach without disturbing the state the \term{SD card}'s boot script already set up.

|Run > Debug Configurations... > GDB OpenOCD Debugging > New Configuration|.

On the \term{Main} tab, set \term{Project} to |hello_world| and \term{C/C++ Application} to |bin/hello_word.elf|.

Leave the build setting on ``Select Automatically'': building this project is |make|'s job, so this never actually triggers a build from \term{Eclipse}.

\begin{figure}[htbp]
    \centering
    \includegraphics[width=\linewidth]{FPGAsteelEclipseDebugConfigMain.png}
    \caption{The \term{Main} tab.}
    \label{fig:fpgasteel-hello-debug-main}
\end{figure}

On the \term{Debugger} tab, enable ``Start OpenOCD locally'', pointing it at the \term{OpenOCD} binary installed in Paragraph~\ref{subsec:fpgasteel-install-paths}, and set \term{Config options} to both config files plus a non-destructive |init_reset| (the board must already be sitting at ``Ready for JTAG'' via the \term{SD card}, Section~\ref{sec:fpgasteel-hello-prerequisites}; this must not reset it):

\begin{shellcode}
-f /home/wheel/projects/FPGAsteel/sw/hello_word/openocd/altera-usb-blaster2.cfg
-f /home/wheel/projects/FPGAsteel/sw/hello_word/openocd/altera_fpgasoc_de.cfg
-c "init; reset init"
\end{shellcode}

Under \term{GDB Client Setup}, point ``Start GDB session'' at the |arm-none-eabi-gdb| from the same toolchain install and add two \term{Commands}: |set mem inaccessible-by-default off| (so \term{GDB} doesn't refuse to read/write memory it hasn't seen mapped) and |set remotetimeout 5| (gives \term{GDB} more time to wait on \term{OpenOCD}'s replies).

\begin{figure}[htbp]
    \centering
    \includegraphics[width=\linewidth]{FPGAsteelEclipseDebugConfigDebugger}
    \caption{The \term{Debugger} tab.}
    \label{fig:fpgasteel-hello-debug-debugger}
\end{figure}

On the \term{Startup} tab, leave ``Initial Reset'' and ``Pre-run/Restart reset'' both unchecked: resetting the \term{HPS} here would lose the \term{DDR}/\term{FPGA}/bridge/cache-\term{MMU} state the \term{SD card} prepared. Check ``Enable \term{Arm} semihosting'' and set both ``Load symbols'' and ``Load executable'' to use the project's own |hello_word.elf|.

\begin{figure}[htbp]
    \centering
    \includegraphics[width=\linewidth]{FPGAsteelEclipseDebugConfigStartup}
    \caption{The \term{Startup} tab.}
    \label{fig:fpgasteel-hello-debug-startup}
\end{figure}

\subsection{Start debugging}
\label{subsec:fpgasteel-hello-start-debugging}

Power on the board and make sure it comes to a stop at the \term{U-Boot} prompt (``Ready for JTAG'', Section~\ref{sec:fpgasteel-hello-prerequisites}): if it doesn't, check the \term{SD card} setup before going any further. Only then press \term{Debug}: \term{OpenOCD} connects over \term{JTAG}, halts the \term{CPU}, loads |hello_word.elf|, stops at the entry point, and is ready for step/breakpoint/run.
